\section{Strict filtered elimination of a degree-one regular symbol}

Let $R$ be a commutative ring, let $I,N\subset R$ be ideals and let
$v\in N$. Suppose multiplication by the degree-one symbol of $v$ is
injective on every piece of the $N$-adic associated graded of $R/I$.
In terms of the entire original ideals, this means that for every
$q\geq0$ and $b\in N^q$,
\[
 vb\in I+N^{q+2}\quad\Longrightarrow\quad b\in I+N^{q+1}.
\]
This is the usual regular-symbol hypothesis; no freeness, geometric
normal-flatness or equality of two filtrations is assumed.

\begin{theorem}[Regular symbols give exact successor division in every finite degree]
\label{hir:regular-symbol-whole-successor-division}
\lean{Hironaka.HironakaStrictFilteredRegularElement.mem_successor_iff_mem_current_of_regular_symbol}
\leanok
For every $q\geq0$ and every $a\in R$,
\[
 va\in I+N^{q+1}\quad\Longleftrightarrow\quad a\in I+N^q.
\]
Thus the entire successor denominator has exactly the expected
principal colon denominator. No separatedness or Noetherian
hypothesis is needed for this finite-degree identity.
\end{theorem}
\begin{proof}
The implication from right to left holds at every degree: write
$a=i+b$ with $i\in I$ and $b\in N^q$. Then $vi\in I$ and
$vb\in N^{q+1}$ because $v\in N$.

For the converse use induction on $q$. At $q=0$ the conclusion is
$a\in I+R=R$. Suppose the converse holds at $q$ and
$va\in I+N^{q+2}$. Since $N^{q+2}\subset N^{q+1}$, the induction
hypothesis writes $a=i+b$ with $i\in I$ and $b\in N^q$.
Subtracting $vi\in I$ from the given inclusion gives
$vb\in I+N^{q+2}$. Injectivity of the degree-$q$ symbol therefore
gives $b\in I+N^{q+1}$. Adding back $i$ gives
$a\in I+N^{q+1}$, completing the induction. The proof uses only
finite powers and whole ideal membership.
\end{proof}

\begin{theorem}[The entire principal image meets the filtration strictly]
\label{hir:regular-symbol-whole-principal-intersection}
\uses{hir:regular-symbol-whole-successor-division}
\lean{Hironaka.HironakaStrictFilteredRegularElement.sup_pow_inf_principal_of_regular_symbol}
\leanok
For every $q\geq0$, the following equality holds for entire ideals:
\[
 (I+N^{q+1})\cap(I+(v))=I+(v)N^q.
\]
In particular, imposing the actual element $v$ introduces precisely
its degree-one graded image in each finite filtration quotient.
\end{theorem}
\begin{proof}
Let $x$ belong to the left side. Write $x=i+va$ with $i\in I$.
Since $x\in I+N^{q+1}$, subtraction of $i$ gives
$va\in I+N^{q+1}$. The preceding theorem gives $a=j+b$ with
$j\in I$ and $b\in N^q$. Hence
$x=(i+vj)+vb\in I+(v)N^q$.

Conversely $I$ is contained in both ideals on the left.
Since $v\in N$, the entire product $(v)N^q$ is contained in
$N^{q+1}$ and is also contained in $(v)$. Therefore its sum with
$I$ belongs to both ideals on the left, proving the reverse inclusion.
This proves equality of whole ideals, including arbitrary original
multipliers. It does not infer that $v$ is injective on the entire
unfiltered quotient: that stronger conclusion requires a separate
intersection-of-all-powers assertion.
\end{proof}
